\section{Hai điểm số đo rò rỉ}
\label{sec:ch15-do-ro-ri}

\index{rò rỉ!vì sao điểm đánh giá không thấy|(}%
Câu mà mục trước dừng lại được chỉ đứng vững nếu rò rỉ là một đại lượng đo được,
vì một bản sao của độ trung thực mà không ai đo được thì cũng chỉ là một phép so
sánh. Đo được thì nó đo được bằng phần chênh mà mục trước để lại chưa có công
thức: chênh giữa lượng thông tin mà cột $\hat c_k$ mang và lượng thông tin mà
khái niệm thật mang. Trước khi viết công thức ấy ra, mục này hỏi vì sao phải
viết, tức những cách đo sẵn có hỏng ở đâu.

\index{rò rỉ!điểm đánh giá khái niệm không thấy}%
Cách đo hiển nhiên nhất là chấm điểm chính bộ mã hóa: độ chính xác trên khái
niệm, điểm F1, hay AUROC. Cách ấy không thấy \tn{rò rỉ}{leakage}, và bài báo cho
thấy điều đó bằng một cặp mô hình chứ không bằng lập luận~\autocite{p30leakage}.
Trên bộ dữ
liệu bảng tổng hợp của nó, một mô hình mềm và một mô hình logit huấn luyện ở
cùng mức $\beta$ có độ chính xác khái niệm $0.993$ và $0.995$, điểm F1 khái niệm
$0.992$ và $0.995$, AUROC khái niệm $0.992$ và $0.995$, rồi độ chính xác, F1 và
AUROC trên nhiệm vụ lần lượt $0.990$ và $0.991$ ở cả ba. Sáu con số, và không
con số nào lệch quá $0.005$. Lấy ngưỡng $0.5$ để đọc một xác suất mềm thành một bit
thì hai mô hình gần như không phân biệt được. Nhưng đem can thiệp lên toàn bộ
khái niệm rồi so với mốc, mô hình thứ nhất mất $0.000$ độ chính xác còn mô hình
thứ hai mất $0.301$.

\index{điểm số can thiệp!định nghĩa và giới hạn}%
Đại lượng vừa dùng cần một cái tên, vì cả mục~\ref{sec:ch15-dung-cu-duoc-kiem}
xoay quanh nó. Gọi nó là \tn{điểm số can thiệp}{intervention score}. Nó là hiệu
của hai độ chính xác nhiệm vụ, lấy theo chiều để nó đo một khoản mất chứ không
đo một khoản được. Số bị trừ là mốc: độ chính xác của một đầu phân loại huấn
luyện riêng trên khái niệm thật, tức mức cao nhất mà một đầu phân loại như thế
với tới được trên bộ dữ liệu ấy. Số trừ là độ chính xác của mô hình đang xét sau
khi đã can thiệp lên tất cả các khái niệm của nó. Phép trừ ấy có lý do: đầu phân
loại có thể đơn giản tới mức không học nổi quan hệ giữa khái niệm và nhãn, và
khi đó độ chính xác thấp mà không phải do rò rỉ, nên lấy mốc ra thì phần còn lại
đúng là phần rò rỉ lấy đi. Với một mô hình cứng huấn luyện thành công thì hai vế
bằng nhau và điểm số can thiệp bằng 0 theo cách dựng; càng rò rỉ thì nó càng
lớn. Bài báo viết nó là $S_{int}$; sách không lấy ký hiệu ấy, vì chữ $S$ đã
thuộc về liên minh đặc trưng của chương~\ref{ch:shap}, và cái tên tiếng Việt đủ
dùng.

\index{điểm số can thiệp!đặc hiệu cao, nhạy thấp}%
Điểm số can thiệp lại không dùng được làm dụng cụ đo chính, và bài báo nói rõ vì
sao. Nó đặc hiệu hoàn hảo mà không nhạy: can thiệp mà kết quả tệ
đi thì chắc chắn có rò rỉ, còn can thiệp mà kết quả tốt lên thì không kết luận
được gì. Nó không tách được hai loại rò rỉ khỏi nhau và không chấm điểm cho từng
khái niệm riêng, nên không chỉ ra chỗ nào trong bảng khái niệm đáng ngờ. Nó đắt,
vì phải chạy lại mô hình sau mỗi bước can thiệp, và đắt nhất đúng ở những bộ dữ
liệu thật có nhiều khái niệm. Và trong các kiến trúc phức tạp hơn nó thôi không
còn là phép đo rò rỉ nữa: bài báo ghi rằng ở lớp mô hình mà
mục~\ref{sec:ch15-nguyen-nhan} đọc, mô hình rò rỉ nhiều hơn lại có thể can thiệp
tốt hơn.

\index{rò rỉ!hai điểm số cũ và chỗ hỏng của chúng}%
Hai điểm số đã được đề xuất trước đó cũng không thay được. Điểm thứ nhất, gọi là
OIS, huấn luyện $\lvert\mathcal{C}\rvert(\lvert\mathcal{C}\rvert - 1)$ mạng phụ,
mỗi mạng học đoán khái niệm thật thứ $l$ từ biểu diễn học được của khái niệm thứ
$k$. Xếp các AUROC thu được thành một ma trận, làm cùng như vậy với khái niệm thật
để có ma trận mốc, rồi lấy chuẩn Frobenius của hiệu hai ma trận, chuẩn hóa về
khoảng từ 0 đến 1. Ý tưởng đúng hướng nhưng ba chỗ hỏng. Nó chỉ nhìn quan hệ
giữa các khái niệm với nhau, nên không thấy loại rò rỉ chính là loại về nhãn. Nó
không tắt trên mô hình cứng: bài báo đo OIS trên mô hình cứng ở tám bộ dữ liệu
và không lần nào giá trị gần 0, mà lại gần bằng giá trị của những mô hình rò rỉ
nặng, nên nó không phân biệt được mô hình diễn giải được với mô hình không. Và
vì định nghĩa của nó chứa một phép huấn luyện, giá trị của nó dao động mạnh giữa
các lần chạy trên cùng một mô hình, đủ mạnh để khoảng tin cậy của hai mô hình
khác hẳn nhau về rò rỉ vẫn chồng lên nhau. Điểm thứ hai, gọi là NIS, mở rộng ý
tưởng ấy sang những tập con khái niệm thay vì từng cặp, và mắc cùng chỗ hỏng thứ
hai: nó cũng không tắt trên mô hình cứng.

\index{rò rỉ!vì sao điểm đánh giá không thấy|)}%
Chỗ chung của cả bốn cách đo trên là chúng đo bằng điểm số dự đoán. Rò rỉ thì
không phải một chuyện về dự đoán mà là một chuyện về lượng thông tin, nên
mục~\ref{sec:ch15-ro-ri} đã phải mượn entropy và thông tin tương hỗ mới nói được
nó là gì. Định nghĩa sau đây chỉ là câu ấy viết thành công thức, và nó cho hai
điểm số mà cả chương sau đó gọi bằng tên viết tắt của chúng: CTL cho rò rỉ khái
niệm-nhiệm vụ và ICL cho rò rỉ liên khái niệm, cả hai theo phụ
lục~\ref{app:viet-tat}.

\index{rò rỉ!hai điểm số thông tin|(}%
\begin{definition}[Hai điểm số rò rỉ]
\label{def:ch15-hai-diem-so}
Với mỗi khái niệm $k$ trong tập khái niệm $\mathcal{C}$, điểm
\tn{rò rỉ khái niệm-nhiệm vụ}{concepts-task leakage} là phần chênh, chặn dưới ở
0, giữa lượng thông tin chuẩn hóa
mà giá trị học được mang về nhãn và lượng thông tin chuẩn hóa mà khái niệm thật
mang về nhãn,
\[
  \mathrm{CTL}_k
  = \max\!\left\{0,\; \frac{I(\hat c_k, y)}{H(y)} - \frac{I(c_k, y)}{H(y)}\right\} .
\]
Với mỗi cặp khái niệm $k \ne l$, điểm \tn{rò rỉ liên khái niệm}{interconcept
leakage} là phần chênh tương ứng giữa hai giá trị học được và giữa hai khái niệm
thật,
\[
  \mathrm{ICL}_{kl}
  = \max\!\left\{0,\;
    \frac{I(\hat c_k, \hat c_l)}{\sqrt{H(\hat c_k)H(\hat c_l)}}
    - \frac{I(c_k, c_l)}{\sqrt{H(c_k)H(c_l)}}\right\} .
\]
Điểm rò rỉ của cả mô hình là trung bình của $\mathrm{CTL}_k$ trên các khái niệm,
và trung bình của $\mathrm{ICL}_{kl}$ trên các cặp~\autocite{p30leakage}.
\end{definition}

Ba chi tiết trong định nghĩa ấy quyết định nó đọc được thế nào. Phép chuẩn hóa bằng entropy khiến
cả hai điểm số nằm trong khoảng từ 0 đến 1, nên so được giữa các bộ dữ liệu có
số lớp khác nhau. Phép chặn dưới ở 0 xử lý trường hợp phần chênh âm, tức giá trị
học được mang \emph{ít} thông tin hơn khái niệm thật: đó là học khái niệm kém
chứ không phải rò rỉ, và gộp hai thứ vào một con số thì con số ấy không đọc được
nữa. Còn $\mathrm{ICL}_{kk}$ bằng 0 theo công thức, vì thông tin tương hỗ của
một biến với chính nó là entropy của nó, nên đường chéo tự triệt tiêu ở cả hai
vế.

Cả hai đại lượng đều là thông tin tương hỗ trên các biến liên tục, thứ không
tính thẳng được từ mẫu hữu hạn. Bài báo ước lượng chúng bằng ước lượng của
Kraskov, Stögbauer và Grassberger, dựa trên thống kê láng giềng gần nhất. Phép
ước lượng ấy là một
lựa chọn tự do của khung, và mục~\ref{sec:ch15-dung-cu-duoc-kiem} sẽ thấy nó đặt
một giới hạn thật lên chỗ khung dùng được.

\index{rò rỉ!hai điểm số thông tin|)}%
Hai điểm số cho hai con số, và so hai mô hình bằng hai con số thì có lúc chúng
không cùng chiều. Bài báo phát biểu sẵn một tiêu chí cho phép kết luận: mô hình
A rò rỉ nhiều hơn mô hình B nếu cả hai điểm số của A đều cao hơn với độ tin cậy
thống kê cao, hoặc nếu một điểm số của A cao hơn với độ tin cậy cao còn điểm số
kia nhận giá trị tương thích thống kê ở hai mô hình. Tiêu chí ấy cố tình dè dặt
và bài báo nói vậy: nó không kết luận gì khi một điểm số cao hơn ở A còn điểm
kia cao hơn ở B, và mọi phát biểu \enquote{rò rỉ nhiều hơn} ở các mục sau đều là
phát biểu theo tiêu chí này chứ không phải theo một phép so tùy ý.
